%-------------------------
% 中文简历（由 main.tex 翻译）
% 请使用 XeLaTeX 编译
%-------------------------

\documentclass[letterpaper,11pt]{article}

\usepackage{latexsym}
\usepackage[empty]{fullpage}
\usepackage{titlesec}
\usepackage{marvosym}
\usepackage[usenames,dvipsnames]{color}
\usepackage{verbatim}
\usepackage{enumitem}
\usepackage[hidelinks]{hyperref}
\usepackage{fancyhdr}
\usepackage{tabularx}
\usepackage{fontspec}
\usepackage{xeCJK}

\setmainfont{Times New Roman}
\setCJKmainfont{Microsoft YaHei}

\pagestyle{fancy}
\fancyhf{}
\fancyfoot{}
\renewcommand{\headrulewidth}{0pt}
\renewcommand{\footrulewidth}{0pt}

% 页边距
\addtolength{\oddsidemargin}{-0.5in}
\addtolength{\evensidemargin}{-0.5in}
\addtolength{\textwidth}{1in}
\addtolength{\topmargin}{-.5in}
\addtolength{\textheight}{1.0in}

\urlstyle{same}
\raggedbottom
\raggedright
\setlength{\tabcolsep}{0in}

% 章节样式
\titleformat{\section}{
  \vspace{-4pt}\raggedright\large
}{}{0em}{}[\color{black}\titlerule \vspace{-5pt}]

% 在支持时提升 PDF 文本可检索性
\ifdefined\pdfgentounicode
  \pdfgentounicode=1
\fi

%-------------------------
% 自定义命令
\newcommand{\resumeItem}[1]{
  \item\small{
    {#1 \vspace{-2pt}}
  }
}

\newcommand{\resumeSubheading}[4]{
  \vspace{-2pt}\item
    \begin{tabular*}{0.97\textwidth}[t]{l@{\extracolsep{\fill}}r}
      \textbf{#1} & #2 \\
      \textit{\small#3} & \textit{\small #4} \\
    \end{tabular*}\vspace{-7pt}
}

\newcommand{\resumeSubSubheading}[2]{
    \item
    \begin{tabular*}{0.97\textwidth}{l@{\extracolsep{\fill}}r}
      \textit{\small#1} & \textit{\small #2} \\
    \end{tabular*}\vspace{-7pt}
}

\newcommand{\resumeProjectHeading}[2]{
    \item
    \begin{tabular*}{0.97\textwidth}{l@{\extracolsep{\fill}}r}
      \small#1 & #2 \\
    \end{tabular*}\vspace{-7pt}
}

\newcommand{\resumeSubItem}[1]{\resumeItem{#1}\vspace{-4pt}}
\renewcommand\labelitemii{$\vcenter{\hbox{\tiny$\bullet$}}$}
\newcommand{\resumeSubHeadingListStart}{\begin{itemize}[leftmargin=0.15in, label={}]}
\newcommand{\resumeSubHeadingListEnd}{\end{itemize}}
\newcommand{\resumeItemListStart}{\begin{itemize}}
\newcommand{\resumeItemListEnd}{\end{itemize}\vspace{-5pt}}

%-------------------------------------------
\begin{document}

%----------个人信息----------
\begin{center}
    \textbf{\Huge 薛睿泽} \\ \vspace{1pt}
    \small \textit{Planning & control & motor_driver} \\
    \href{mailto:xrz836884211@gmail.com}{\underline{xrz836884211@gmail.com}} $|$
    18640745980 $|$
    深圳
\end{center}

%----------求职意向----------
\section{求职意向}
\small{自动化工程师 / 机器人控制算法工程师 / 嵌入式算法工程师}

%----------个人优势----------
\section{个人优势}
\small{
机器人与控制工程师，专注于云台控制、嵌入式电机验证及软硬件集成。具备基于系统辨识的电机测试、MCU 端机器学习诊断、可靠性标准制定及自动化验证流程开发经验，熟悉 Python、C++、MATLAB/Simulink、GDB-JLink、UART 及实验室仪器。研究方向涵盖移动机械臂规划、碰撞规避和尾座式无人机控制。
}

%----------核心技能----------
\section{核心技能}
 \begin{itemize}[leftmargin=0.15in, label={}]
    \small{\item{
     \textbf{机器人与控制}{：电机控制、PID/LQR/MPC、LADRC/ESO、云台传感器融合、轨迹生成} \\
     \textbf{自动化与验证}{：电机及云台测试方案设计与实施、FRF 分析、模态辨识、信号采集、数据分析} \\
     \textbf{软件开发}{：Python、C++、C、MATLAB/Simulink、Git、Linux} \\
     \textbf{嵌入式与硬件}{：MCU 开发、ADC/DMA、PWM、电机驱动、传感器、示波器、数据采集} \\
     \textbf{路径规划与机器学习}{：路径规划、碰撞规避、机器学习} \\
     \textbf{工程设计与制造}{：Fusion 360、SolidWorks、3D 打印}
    }}
 \end{itemize}

%----------工作经历----------
\section{工作经历}
  \resumeSubHeadingListStart
    \resumeSubheading
      {云台控制算法工程师}{2024.07 -- 至今}
      {Insta360}{深圳}
      \resumeItemListStart
        \resumeItem{\normalsize{\textbf{从 0 到 1 建设云台电机测试体系}} \\
        \small{开发基于系统辨识的电机测试方法，定量表征齿槽转矩、摩擦、相电阻、电感及 Hall 传感器健康状态。搭建服务端自动化数据采集、分析与报告生成流程，提高电机量产品质管理中的供应商验证、失效分析及设计决策效率。}}
        \resumeItem{\normalsize{\textbf{面向云台保护的 L-SVM 碰撞检测}} \\
        \small{设计可部署于 MCU 的 L-SVM 碰撞检测器，识别用户操作引发、可能造成电机结构变形及控制性能下降的云台碰撞事件。提出将用户侧校准作为核心产品功能，提升产品鲁棒性，并形成相较竞品云台系统的差异化优势。}}
        \resumeItem{\normalsize{\textbf{电机可靠性测试标准与云台健康检查}} \\
        \small{负责云台项目电机可靠性测试标准的制定，并定义控制侧与电机侧的健康检查方法。通过可靠性问题诊断及电机侧风险验证，持续支持云台产品开发并维护控制性能健康标准。}}
        \resumeItem{\normalsize{\textbf{代码开发与 Jira 验证自动化 Agent}} \\
        \small{开发自动化 Agent 工作流：通过 Lark CLI 读取需求文档，在固件仓库内实现功能，并调用自研 GDB-JLink MCP、UART CLI、数字电源 MCP 及 SoC 相关 CLI 验证功能或缺陷修复。实现验证证据自动收集，并将结果回传至 Lark 或 Jira。}}
        \resumeItem{\normalsize{\textbf{云台基于频域鲁棒优化的控制器生成}} \\
        \small{研究基于频域鲁棒优化的控制器自动生成方法，在简化控制器结构的同时，使性能接近现有人工调参方案。推导闭式数学解，并在 MATLAB 中通过频率响应、稳定裕度、跟踪及抗扰指标验证方法有效性。}}
      \resumeItemListEnd
  \resumeSubHeadingListEnd

%----------兼职科研经历----------
\section{兼职科研经历}
  \resumeSubHeadingListStart
    \resumeSubheading
      {兼职研究员}{2024.12 -- 至今}
      {移动机械臂：面向随机形状末端负载的无碰撞轨迹生成}{RAL 投稿，第三作者}
      \resumeItemListStart
        \resumeItem{\normalsize{\textbf{基于 Hybrid A* 与 RRT* 的前端路径采样}} \\
        \small{开发采样式规划框架：使用基于 A* 的方法为两轮移动底盘生成路径样本，并在六轴机械臂关节空间内使用 RRT* 进行规划，以实现复杂障碍场景下的无碰撞运动。}}
        \resumeItem{\normalsize{\textbf{面向真实实验的工程实现}} \\
        \small{使用 Python、C++ 和 MATLAB 实现规划流程，集成机器人模型、碰撞检测、轨迹导出及实验验证指标，包括规划时间、跟踪误差、安全裕度和成功率。}}
      \resumeItemListEnd
  \resumeSubHeadingListEnd

%----------其他相关经历----------
\section{其他相关经历}
  \resumeSubHeadingListStart
    \resumeSubheading
      {机器人研发实习生}{2023.12 -- 2024.06}
      {东莞市本末科技（DiaBot Technology）}{东莞}
      \resumeItemListStart
        \resumeItem{\normalsize{\textbf{轮腿机器人与机械臂抓取控制}} \\
        \small{适配合作方初创公司的自研机械臂并评估轻量化替代方案。设计基于几何法的逆运动学求解器，实现头部稳定功能；将正逆运动学、插值及串级控制算法移植至轮腿机器人运动控制器，并去除 ROS 依赖。交付轮腿机器人与机械臂演示系统，曾面向华为访客及李教授创业孵化基地展示。}}
        \resumeItem{\normalsize{\textbf{Tita 机器人 Core-XY 换电站演示系统}} \\
        \small{与机械工程师协作，从 0 到 1 共同设计换电站原型。采用 Core-XY 结构和步进电机构建 XYZ 运动控制框架，实现机器人进站后自动更换电池。参与专利 CN118559741A，目前处于实质审查阶段。}}
        \resumeItem{\normalsize{\textbf{Tita 轮腿机器人跳跃冲击测量}} \\
        \small{基于高频 IMU、UDP 和 ESP32 搭建冲击数据采集设备，并使用 MATLAB App Designer 开发上位机，实现 IMU 参数标定及基于 EKF 的姿态估计。测试设备已交付机械团队并持续在内部使用。参与专利 CN118330258A，目前处于实质审查阶段。}}
      \resumeItemListEnd

    \resumeSubheading
      {尾座式垂直起降无人机设计与控制}{2022.09 -- 2023.08}
      {MarsLab 毕业设计}{香港大学}
      \resumeItemListStart
        \resumeItem{\normalsize{\textbf{高效垂直起降无人机控制}} \\
        \small{推导尾座式无人机的运动学与动力学模型，在仿真环境中设计串级 PID 控制器，快速验证控制逻辑与导航行为。负责动力学仿真及控制器设计，完成定点导航仿真，并支持基于 Vicon 与 PX4 硬件平台的室内实验。}}
      \resumeItemListEnd
  \resumeSubHeadingListEnd

%----------教育背景----------
\section{教育背景}
  \resumeSubHeadingListStart
    \resumeSubheading
      {香港大学（HKU）}{香港}
      {机械工程理学硕士，MarsLab}{2022.09 -- 2024.07}
    \resumeSubheading
      {匹兹堡大学}{美国宾夕法尼亚州匹兹堡}
      {机械工程理学学士}{2021.09 -- 2022.05}
    \resumeSubheading
      {四川大学}{成都}
      {机械设计制造及其自动化工学学士}{2018.09 -- 2021.07}
  \resumeSubHeadingListEnd
  \begin{itemize}[leftmargin=0.15in, label={}]
    \small{\item{
      \textbf{相关课程}{：机电一体化系统工程、无人机设计控制与导航、现代机器人学、机器学习与模式识别、机电一体化、自动控制、机器人学与控制}
    }}
  \end{itemize}

%----------论文与奖项----------
\section{论文与奖项}
  \begin{itemize}[leftmargin=0.15in, label={}]
    \small{\item{
      \textbf{Swarm-LIO2: Decentralized Efficient LiDAR-Inertial Odometry for Aerial Swarm Systems} \\
      \textbf{Autonomous Tail-Sitter Flights in Unknown Environments} \\
      \textbf{AI 能力认证奖}{：Insta360 嵌入式软件部}
    }}
 \end{itemize}

%----------其他信息----------
\section{其他信息}
  \begin{itemize}[leftmargin=0.15in, label={}]
    \small{\item{
      \textbf{语言能力}{：英语口语流利，可适应全英文工程协作环境。} \\
      \textbf{招聘支持}{：通过专业人脉协助工程团队开展兼职技术招聘，搜寻并评估高质量候选人。} \\
      \textbf{个人兴趣}{：养有一只狗和一只猫。}
    }}
 \end{itemize}

%-------------------------------------------
\end{document}
